\Qhead{2.6}{Волноводы, световоды, линзовые линии и открытые резонаторы}{Электромагнитные волны в металлических волноводах. Диэлектрические волноводы, световоды. Линзовые линии и открытые резонаторы. Гауссовские пучки.}
\Subfile{2.6.1}{sections/02_waves/q06/s01}
\Subfile{2.6.2}{sections/02_waves/q06/s02}
\Subfile{2.6.3}{sections/02_waves/q06/s03}

\recap{Поперечные граничные условия создают моды и отсечки. Металлический волновод удерживает поле проводящими стенками, диэлектрический — контрастом показателя; открытый резонатор — самовоспроизведением дифрагирующего пучка.}
\board{Поперечная собственная задача; TE/TM/TEM; условие $n_2k_0<\beta<n_1k_0$; гауссов пучок и $z_R$.}
\oral{Почему TEM-мода возможна в коаксиальной линии, но отсутствует в пустом прямоугольном волноводе?}
\bait{Можно оставить фазу Гуи: дополнительный фазовый набег фокусирующегося пучка сдвигает частоты поперечных мод резонатора.}
\sources{Цысарь~\cite{tsysar1}, гауссов пучок: PDF-с.~44--45; открытые резонаторы: с.~46--48; Белокопытов~\cite{belokopytov1}, разделы о направляемых волнах; Митрофанов~\cite{mitrofanov}, лекция~13. Граничные условия и гауссова нормировка выписаны явно.}
